% ===========================================================================
% 01 - Einordnung: Agenda, Ziel und Traceability-Ansatz
% ===========================================================================

\begin{frame}{Agenda}
  \begin{columns}[T,onlytextwidth]
    \begin{column}{0.48\textwidth}
      \textbf{Teil 1 -- Einordnung}
      \begin{itemize}
        \item Ziel der Präsentation
        \item Traceability-Ansatz
      \end{itemize}
      \vspace{0.4cm}
      \textbf{Teil 2 -- Architektur}
      \begin{itemize}
        \item Service-Landschaft
        \item Choreographie
      \end{itemize}
    \end{column}
    \begin{column}{0.48\textwidth}
      \textbf{Teil 3 -- Technische Anforderungen}
      \begin{itemize}
        \item TA-01 bis TA-05 im Code
      \end{itemize}
      \vspace{0.4cm}
      \textbf{Teil 4 -- Architekturentscheidungen}
      \begin{itemize}
        \item ADR-01 bis ADR-06 im Code
      \end{itemize}
      \vspace{0.4cm}
      \textbf{Teil 5 -- Demo \& Fazit}
    \end{column}
  \end{columns}

  \vfill
  \begin{block}{Leitfrage}
    \small Welche Anforderung ist \textbf{wo} im Code nachweisbar -- und
    \textbf{warum} wurde sie so umgesetzt?
  \end{block}
\end{frame}

\begin{frame}{Ziel: Nachweis statt Behauptung}
  \begin{columns}[T,onlytextwidth]
    \begin{column}{0.54\textwidth}
      \textbf{Problem}
      \begin{itemize}
        \item Architekturdokumente beschreiben Anforderungen.
        \item Ohne Code-Bezug bleibt offen, ob sie umgesetzt sind.
      \end{itemize}

      \vspace{0.3cm}
      \textbf{Lösung: Traceability}
      \begin{itemize}
        \item Jede Anforderung $\rightarrow$ konkrete Datei + Symbol.
        \item Jedes Snippet $\rightarrow$ kurze Erläuterung.
        \item Nachvollziehbar in der mündlichen Prüfung.
      \end{itemize}
    \end{column}
    \begin{column}{0.42\textwidth}
      \begin{block}{Kennungen}
        \small
        \begin{tabular}{@{}ll@{}}
          \tabadge{TA-xx} & Technische Anforderung \\
          \adrbadge{ADR-xx} & Entwurfsentscheidung \\
          \frbadge{FR-xx} & Funktionale Anforderung \\
          \qrbadge{QR-xx} & Qualitätsanforderung \\
        \end{tabular}
      \end{block}
      \vspace{0.2cm}
      \begin{block}{Quelle}
        \small \code{docs/code-traceability.md}
      \end{block}
    \end{column}
  \end{columns}
\end{frame}

\begin{frame}{Traceability-Matrix (Auszug)}
  \small
  \begin{tabularx}{\textwidth}{L{1.6cm}L{4.6cm}Y}
    \toprule
    \textbf{Kennung} & \textbf{Anforderung} & \textbf{Zentraler Code-Nachweis} \\
    \midrule
    TA-01 & Spring Boot & \code{@SpringBootApplication} in 6 Service-Klassen \\
    TA-02 & Docker & 6 $\times$ \code{Dockerfile} (Multi-Stage) \\
    TA-03 & Docker Compose & \code{docker-compose.yml} (6 Services, Volumes) \\
    TA-04 & Resilience4j & \code{@CircuitBreaker} + \code{fallback} + \code{application.yml} \\
    TA-05 & Kafka (COULD) & bewusst nicht umgesetzt (ADR-02) \\
    \midrule
    ADR-01 & Fachlicher Schnitt & Modulstruktur \code{api/application/domain/infrastructure} \\
    ADR-03 & Choreographie & \code{AnalysisWorker} $\rightarrow$ \code{NextServiceClient} \\
    ADR-06 & Keine Shared Persistence & \code{ConfigurationSnapshot} + stateless Services \\
    \bottomrule
  \end{tabularx}

  \vfill
  \begin{center}
    \small\color{codegray}
    Vollständige Matrix: \code{docs/code-traceability.md}
  \end{center}
\end{frame}
